\begin{table}[htbp]
\centering
\footnotesize
\setlength{\tabcolsep}{4pt}
\caption{两步反演在各模拟设定下给出的隐含环境方差 $v_R^{\mathrm{impl}}$，以及各估计值的 95\% 置信区间。}
\label{tab:s8}
\begin{tabular}{lllrrll}
\toprule
设定 & 由 $R_A$ 得到的 $v_R^{\mathrm{impl}}$ & 由 $G$ 得到的 $v_R^{\mathrm{impl}}$ & $\kappa_{R_A}$ & $\kappa_G$ & 比值 $R_A/G$ & 一致 \\
\midrule
数据匹配设计（M0） & 0.040 [0.026, 0.056] & 0.043 [0.025, 0.064] & 0.516 & 2.729 & 0.93 [0.78, 1.18] & 是 \\
无全国冲击 & 0.039 [0.026, 0.055] & 0.042 [0.024, 0.062] & 0.523 & 2.775 & 0.94 [0.79, 1.18] & 是 \\
全国份额 0.65 & 0.040 [0.026, 0.056] & 0.043 [0.025, 0.065] & 0.510 & 2.670 & 0.92 [0.77, 1.16] & 是 \\
AR(1)，$\phi=0.5$ & 0.027 [0.017, 0.040] & 0.025 [0.013, 0.039] & 0.780 & 4.699 & 1.08 [0.88, 1.43] & 是 \\
真实代际间隔均值 2.2 天 & 0.013 [0.008, 0.019] & 0.014 [0.008, 0.021] & 1.287 & 6.991 & 0.94 [0.77, 1.20] & 是 \\
真实代际间隔均值 4.3 天 & 0.058 [0.038, 0.081] & 0.063 [0.037, 0.095] & 0.380 & 1.977 & 0.92 [0.77, 1.16] & 是 \\
延迟均值 3.5 天 & 0.038 [0.025, 0.053] & 0.041 [0.024, 0.061] & 0.535 & 2.790 & 0.93 [0.77, 1.16] & 是 \\
延迟均值 8 天 & 0.046 [0.030, 0.064] & 0.047 [0.027, 0.069] & 0.454 & 2.483 & 0.98 [0.82, 1.22] & 是 \\
$k=0.1$ & 0.040 [0.026, 0.055] & 0.042 [0.025, 0.062] & 0.517 & 2.740 & 0.94 [0.79, 1.18] & 是 \\
报告计数同步冲击 & 0.019 [0.006, 0.035] & 0.029 [0.012, 0.048] & 1.041 & 4.553 & 0.67 [0.39, 0.90] & 否 \\
通用设计（衔接） & 0.043 [0.029, 0.059] & 0.044 [0.027, 0.063] & 0.507 & 2.734 & 0.97 [0.83, 1.18] & -- \\
\bottomrule
\end{tabular}
\par\smallskip
\begin{minipage}{0.95\linewidth}\scriptsize 各设定只改变数据匹配设计中的一个要素（第~\ref{sec:meth-inversion}~节），分析者所用代际间隔的均值始终为 3.6 天。区间合并了两个实证截距的 2000 次联合双向自助抽样，以及校准线按其蒙特卡罗抽样分布的抽取。$\kappa=\E(\hat{a})/(v_R/\bar{R}^2)$，取最接近隐含方差的 $v_R$ 格点值。“一致”表示两个隐含方差之差的 95\% 区间包含零；第一套模拟的通用设计作为衔接设定列出，不计入十种设定。\end{minipage}
\end{table}
